\section{Introduction}
\IEEEPARstart{B}{us} control must reconcile passenger service with vehicle circulation. Holding a bus can reduce an excessive gap behind it, but the same intervention increases the time spent on board and may postpone the vehicle's next terminal departure. Adjusting a terminal release can redistribute service toward changing demand, but its intended benefit depends on whether downstream controls preserve the revised plan. These interactions motivate a joint treatment of timetabling and holding under a fixed fleet and service commitment.

Headway regulation provides a well-established means of controlling bunching. Daganzo's adaptive holding formulation uses real-time headway information to reduce the slack needed for reliable service~\cite{daganzo2009headway}. Cooperation between successive vehicles also addresses the propagation of large gaps~\cite{daganzo2011cooperation}. Learning-based approaches extend the information and temporal scope of bus control: Wang and Sun study fleet coordination through multi-agent holding~\cite{wang2020dynamic}, and their asynchronous formulation explicitly represents the fact that vehicles request control at different times~\cite{wang2021asynchronous}. These approaches motivate the use of operational state and future consequences rather than a purely local spacing error.

Joint control introduces an additional difficulty. The two layers can share observations while still disagreeing about resources. A slow policy may publish an earlier departure on the assumption that a returning vehicle will be available. A fast controller may subsequently allocate holding to that vehicle because its local prediction improves passenger cost or regularity. Conversely, a budget intended to limit total holding can be ineffective if it is recreated at each control point or if the prediction applies a different action projection from execution. In either case, a nominally hierarchical policy lacks a consistent physical meaning.

Frequency information offers a possible organizing principle. Slowly varying demand summaries can inform limited timetable revisions, whereas local control reacts to current vehicle states and recent service observations. However, three different statements must be separated: observations contain multiple temporal components; actions operate at different frequencies; and assigning particular components to particular actions improves control. The first two do not establish the third. In particular, automatic passenger counts (APC) measure boarding events. Their high-frequency variation can reflect vehicle service timing as well as changes in passenger arrivals. Treating that variation as an immediate demand correction can be misleading.

We present FreqDuet as an implemented, frequency-informed coordination architecture and examine it through bounded development experiments. A slow policy selects between maintaining a persistent holding budget and jointly applying that budget with a demand-informed timetable revision. A fast controller enumerates the legal holding actions at the current event, evaluates passenger cost under a common continuation, and respects the published plan and remaining budget. The learned decision is a permission to apply a defined timetable law; it is not an unrestricted generator of an entire daily schedule. This restriction makes the interaction inspectable and permits matched comparisons with manual control.

The contributions are threefold. First, we specify a shared interface for executable timetable revisions and cumulative holding allowances, with operating timescales and budget quantities that can be inspected by a dispatcher. Second, paired comparisons distinguish upper-policy choice at matched budget retention from the service--regularity effects of changing retention, while matched retraining evaluates harmonic representation separately. Third, a same-state capture traces a timetable revision to an actual holding change, while closed-loop diagnostics explain why better continuous predictions need not produce better discrete responses.

The evaluated system keeps the fleet and trip inventory fixed. Its operational configuration concerns release adjustments and holding time allowances, not additional vehicle or driver allocation. The main experiment compares learned and strong manual joint control under two matched retention settings; a second experiment changes retention within each controller, and a third compares harmonic bases under matched retraining. This separates a configuration benefit from a learning increment. Section~\ref{sec:limitations} collects the evidence boundaries.

\section{Related Work}
\subsection{Holding, Prediction, and Fleet Coordination}
Adaptive headway control relates interventions to the current spacing of vehicles~\cite{daganzo2009headway}. Cooperative control extends this view to interactions with both preceding and following buses~\cite{daganzo2011cooperation}. These ideas are relevant to the native bilateral rule retained as a continuation policy in this study. They also show why a reference rule should remain available when a predictive controller lacks adequate support.

Prediction introduces a distinction between a desirable action under a forecast and its value in physical operation. The comparison of holding methods with and without real-time predictions by Berrebi \emph{et al.}~\cite{berrebi2018comparing} is directly relevant to this distinction. Here, we focus on a more specific interface: future holding decisions, terminal returns, and public departure obligations must be represented on the same clock. Our diagnostics therefore compare identical windows and distinguish a timing forecast from the discrete command selected using that forecast.

Multi-agent reinforcement learning can coordinate holding across a fleet~\cite{wang2020dynamic}. The asynchronous actor--critic framework of Wang and Sun~\cite{wang2021asynchronous} addresses non-simultaneous decision events. Our slow decision is instead a system-level choice between two resource-consistent publication modes, while the fast layer uses explicit enumeration. Hierarchical multi-agent bus control is also an established research direction~\cite{yu2024hierarchical}. The contribution sought here is the connection between timetable commitments and vehicle-level holding, not a new general reinforcement-learning optimizer.

\subsection{Temporal Structure in Learning and Control}
Action frequency is itself a design variable. Lee \emph{et al.} formalize reinforcement learning with action components operating at different frequencies~\cite{lee2020multiple}. Metelli \emph{et al.} analyze control-frequency adaptation through action persistence~\cite{metelli2020persistence}. In a bus system, terminal planning and stop-level holding naturally possess different opportunities for intervention. A committed release remains relevant across many subsequent events, and a holding budget persists until the corresponding trip returns.

FreqDuet additionally uses a causal harmonic representation of demand-related observations. This signal representation should not be conflated with action persistence. The present architecture uses low-frequency summaries in the timetable law and frequency summaries in the slow policy state; it does not inject a high-frequency residual directly into the fast holding command. The distinction is important when interpreting the method name and designing an ablation. Removing frequency information requires changing both the plan law and the policy inputs while maintaining the same operational opportunities.

The upper policy uses the maximum-entropy actor--critic framework~\cite{haarnoja2018sac} with a finite action library, following the discrete-action formulation~\cite{christodoulou2019discrete}. We retain a fixed training budget and compare the policy against its own critics only as a diagnostic. Agreement with a critic is not used as a substitute for a same-state traffic comparison.
